\section{Introduction}
\label{sec:introduction}

Soybean productivity maps at fine spatial resolution would support planning, insurance, and risk assessment, but field-level yield labels are rarely available at scale in Brazil.
What is available instead are official municipal statistics from the Brazilian Institute of Geography and Statistics (IBGE) Municipal Agricultural Production survey (PAM), together with dense Sentinel-2 time series that resolve spectral variation \emph{within} municipalities.
This paper treats that gap as a \emph{downscaling} task: learn pixel-level productivity contributions from municipal labels alone, then read out both municipal aggregates and intramunicipal maps.

Deep Transformers have been shown to extract informative representations from Sentinel-2 time series for crop mapping~\citep{xu2024sitsmoco}.
We adapt the Spectral--Temporal Network (STNet) from that line of research to continuous productivity regression (tonnes per hectare).
Each soybean pixel is scored by the network; municipal predictions are the mean of pixel scores and are supervised by PAM labels, using Sentinel-2 reflectance and Xavier climate covariates over Paran\'{a} for harvest years 2020--2024.

Our contributions are a municipal-to-pixel downscaling scheme that trains per-pixel productivity predictors under mean aggregation to IBGE PAM labels; a spectral--temporal Transformer for Paran\'{a} soybean that combines reflectance with Xavier climate covariates, with ablations that remove climate inputs or change temporal pooling; and a held-out evaluation on the 2021 growing season, including early-season truncations and qualitative intramunicipal maps.
